% 核心结论：求导：乘以角频率并超前九十度
% 图原型：纯几何／对比；按本题相位与电路拓扑绘制。
\documentclass[tikz,border=7pt]{standalone}
\usepackage{fontspec,xeCJK,amsmath,amssymb}
\IfFontExistsTF{Times New Roman}{\setmainfont{Times New Roman}}{\setmainfont{TeX Gyre Termes}}
\IfFontExistsTF{Heiti SC}{\setCJKmainfont[AutoFakeBold=1.5]{Heiti SC}}{\setCJKmainfont[AutoFakeBold=1.5]{FandolHei-Regular}}
\usetikzlibrary{arrows.meta,calc,angles,quotes}
\definecolor{ink}{HTML}{243745}
\definecolor{blue}{HTML}{275DAB}
\definecolor{green}{HTML}{177C63}
\definecolor{red}{HTML}{BD4935}
\definecolor{purple}{HTML}{715393}
\begin{document}
\begin{tikzpicture}[font=\fontsize{9.5}{12}\selectfont,text=ink,
 v/.style={line width=1.1pt,-{Stealth[length=2.4mm,width=1.5mm]}},
 axis/.style={ink!65,line width=.45pt,-{Stealth[length=1.7mm,width=1.1mm]}},
 aux/.style={ink!40,densely dashed,line width=.5pt},
 lab/.style={fill=white,inner sep=2pt},
 card/.style={draw=blue!25,fill=blue!3,rounded corners=3pt,text width=6cm,align=left,inner sep=10pt}]
\path[use as bounding box] (-8,-3.35) rectangle (8,3.35);
\node[font=\fontsize{11}{14}\selectfont\bfseries] at (0.00000,3.04000) {求导：乘以角频率并超前九十度};
\begin{scope}[shift={(-4,0)}]
\draw[axis] (-2.50000,0.00000) -- (2.60000,0.00000);
\draw[axis] (0.00000,-2.05000) -- (0.00000,2.25000);
\node[] at (2.58000,-0.26000) {$\operatorname{Re}$};
\node[] at (-0.33000,2.20000) {$\operatorname{Im}$};
\node[] at (-0.18000,-0.20000) {$O$};
\draw[v,blue] (0.00000,0.00000) -- (1.69145,-0.61564);
\node[lab,text=blue,anchor=center] at (0.94572,-0.65782) {$\widetilde A$};
\draw[v,red] (0.00000,0.00000) -- (0.71824,1.97335);
\node[lab,text=red,anchor=east] at (-0.40000,1.40000) {$j\omega\widetilde A$};
\draw[green] (.9,-.32757) arc[start angle=-20,end angle=70,radius=.95776];
\node[lab,text=green] at (1.45000,0.75000) {$90^\circ$};
\end{scope}
\begin{scope}[shift={(4,0)}]
\node[card] at (0.00000,0.80000) {$\dfrac{\mathrm da}{\mathrm dt}\ \longleftrightarrow\ j\omega\widetilde A$\\[5pt]模变为 $\omega A$，初相增加 $90^\circ$。};
\node[card] at (0.00000,-1.00000) {积分对应 $1/(j\omega)$，\\只描述同频稳态分量。};
\end{scope}
\node[text width=15cm,align=center] at (0.00000,-3.00000) {原量与导数的量纲不同，图中分用尺度；不能由两箭长比较数值。};
\end{tikzpicture}
\end{document}
